% =====================================================================
% Section 3. Numerical method
% Source: docs/derivation/B_TV_IMEX.md (B0-B11, theorems in B9.5),
%         docs/derivation/E_code_path.md (algorithm, settings)
% =====================================================================
\section{Numerical method}
\label{sec:method}

% One paragraph: overview. Transport--velocity (TV) splitting; explicit
% material-speed Rusanov for transport; one implicit pressure equation that
% contains the acoustic part and all relaxation through a single local closure;
% the rates mu, theta, nu are parameters of that closure.

\subsection{Design requirements}
\label{sec:requirements}
% TODO: table R1-R7; every discrete choice below is tied to one of them.
%  R1 conservation of m_k, rho u, E;
%  R2 uniform (p,u) preserved for arbitrary alpha_k, rho_k (well-balanced contacts);
%  R3 alpha_k, m_k, eps_k transported by the same face operator (temperature consistency);
%  R4 the same Delta V_k in volume change and pressure work (discrete Gibbs relation);
%  R5 the implicit pressure equation sees the relaxed stiffness;
%  R6 mixture internal energy consistent with the conserved total energy;
%  R7 one closure, rates as parameters, both limits exact.

\subsection{Transport--velocity splitting}
\label{sec:tv_split}
% TODO: Toro--Vazquez-Cendon splitting; incremental flux form
%   X u^{n+1} = X u^e (explicit) + X (u^{n+1} - u^e) (implicit increment);
% table: explicit / implicit part of each equation;
% eigenvalues of the explicit subsystem are all u -> dissipation speed = material speed;
% mass also carries the increment (temperature consistency, R3).

\subsection{Volume-transport form and path-conservative Rusanov}
\label{sec:volume_transport}
% TODO: linear path, \bar u_f = (u_L+u_R)/2; identity
%   \bar u \Delta\alpha = \Delta(\alpha u) - \bar\alpha \Delta u;
% conservative volume flux + volume strain -\bar\alpha_f \Delta u_f;
% the strain is moved to the implicit closure and paired with the pressure work (R4).

\subsection{Explicit operator}
\label{sec:explicit}
% TODO: MUSCL reconstruction (shared limiter for volume-transported quantities,
% velocity separately); Rusanov flux with material speed
\begin{equation}
F^e_{i+1/2}[X]=\tfrac12\big(X^Lu^L+X^Ru^R\big)-\tfrac12a_f\big(X^R-X^L\big),\qquad a_f\ge\max\big(|u^L_f|,|u^R_f|\big).
\label{eq:rusanov}
\end{equation}
% local or global a_f (per test); positivity:
% Positivity of this operator is part (i) of Theorem~\ref{thm:positivity} (Sec. 3.5.2).
% Remark: any Sweby-TVD limiter works (B9.5 Note 15); minmod is a choice.

\subsection{Implicit operator: face velocity, face pressure and incremental fluxes}
\label{sec:implicit_faces}

\subsubsection{Face velocity and face pressure}
% TODO: momentum-consistent pressure acceleration (half-cell momentum balance),
% density-weighted face pressure; acoustic complement:
\begin{equation}
\begin{aligned}
u_f&=\bar u^*_f-\frac{\tau(P_{i+1}-P_i)}{\bar\rho_f\Delta x}-\varphi_f\,\psi^P_f\,\frac{P_{i+1}-P_i}{Z_i+Z_{i+1}},\\
p_f&=\frac{\rho^*_{i+1}P_i+\rho^*_iP_{i+1}}{\rho^*_i+\rho^*_{i+1}}-\varphi_f\,\psi^u_f\,\frac{Z_iZ_{i+1}}{Z_i+Z_{i+1}}\,D^*_f,
\end{aligned}\qquad
\varphi_f=\Big(1-\frac{\tau\min(c_i,c_{i+1})}{\Delta x}\Big)^+ .
\label{eq:faces}
\end{equation}
% Derivations in the text, not theorems:
%  - weight phi = (1 - nu)^+ from the per-stage acoustic viscosity 1/2 c dx (nu + phi)
%    (B9.5 Thm 6; short von Neumann argument, details in an appendix);
%  - psi^P coupled with the shared limiter, psi^u from u alone: Jameson's LED criterion (B9.5 Thm 7);
%  - frozen sound speed is subcharacteristic, rho c_f >= rho c_eq (remark, cite Flatten--Lund 2011).
% Low-Mach behaviour of the complement (phi = 0 when c tau >= dx).

\subsubsection{Incremental fluxes and transported states}
% TODO: increment I_f = X^L_f w^L + X^R_f w^R with the reconstructed face values
% of the explicit operator as donors (upwind by u_f);
% alpha_2 face value = 1 - alpha_1 face value; sum of face volumes = 1.
\begin{theorem}[Positivity]\label{thm:positivity}
% TODO (B9.5 Thms 4 and 5 merged):
% (i) explicit operator: a_f >= |u^{L,R}_f| and a r <= 1/2;
% (ii) explicit flux + implicit increment with face-value donors: in addition
%      r (a_f + |u_f|) <= 1/2 -> every phase is a non-negative combination of
%      admissible face states, independent of the content ratio.
\end{theorem}

\subsection{Unified implicit closure}
\label{sec:closure}

\subsubsection{Discrete Gibbs relation}
% TODO: backward Euler of the implicit part; volume strain shared with frozen
% fractions omega_k; relaxation terms with chi_k = +-1:
\begin{equation}
\begin{aligned}
\alpha_k&=\alpha^{\rm tr}_k+\omega_ks+\chi_kR,\qquad R=\tau\big(\mu(p_1-p_2)+\mathcal M/\rho_I\big),\\
m_k&=m^{\rm tr}_k+\chi_k\tau\mathcal M,\\
\varepsilon^{\rm nc}_k&=\varepsilon^{\rm tr}_k-\omega_kp_k\,s-\chi_k\,\tau p_I\mu(p_1-p_2)+\chi_k\tau\big(\mathcal Q+e_I\mathcal M\big).
\end{aligned}
\label{eq:gibbs}
\end{equation}

\subsubsection{Shock relations of the model}
\label{sec:shock_relations}
% TODO: the phase energy equations are non-conservative; across a shock the
% conservation laws fix only the mixture Hugoniot e - e_0 + \bar p (v - v_0) = 0.
% The partition of dissipation between phases is set by the shock structure.
% Closure adopted: SPB (II.5) -- each phase on its own Hugoniot,
%   e_k - e_{k,0} + 1/2 (p_0 + p_1)(v_k - v_{k,0}) = 0,
% which is what the energy-defect partition below reproduces.
% Lax admissibility: G_k > 0 -> admissible shocks are compressive (Bethe--Weyl),
% used as the compression criterion of the shock budget (remark).

\subsubsection{Mixture-energy consistency and the energy-defect partition}
% TODO: defect delta = (rho e)^C - sum eps^tr + s sum omega_k p_k;
% eps_k = eps^nc_k + Lambda_k(delta), sum Lambda_k = delta;
% (a) the partition fixes the limit shock relations;
% (b) physical dissipation: share of the phase Hugoniot--isentrope energy gap Q_k (II.5);
% (c) local shock strength from the compact +-1 stencil, compression criterion;
% (d) discretisation error: split by the admissible energy E_k;
% (e) continuity in delta, piecewise affine Lambda_k.
\begin{theorem}[Shock share]\label{thm:shock_share}
% TODO (B9.5 Thms 8 and 9 merged):
% (i) weak shocks: the (II.5)-consistent share is unique, lambda_k ~ alpha_k G_k / K_k^2;
% (ii) no continuous local rule is (II.5)-consistent for all shock strengths
%      (counterexample from the shock-profile ODE) -> the budget needs non-local data.
% Proof of (ii) in an appendix.
\end{theorem}
\begin{proposition}[Error split]\label{prop:error_split}
% TODO (B9.5 Thm 12): lambda^E_k = E_k / sum E is the unique delta-linear split
% keeping both phases admissible for all |delta| < sum E.
% Remark after it: variational characterisation (min sum Lambda_k^2 / E_k, B9.5 Prop 13).
\end{proposition}

\subsubsection{Closures for each set of rates}
% TODO: one subsection-level paragraph each, residual and its monotonicity:
%  mu = inf, theta = nu = 0       : p closure (algebraic, NASG; 1D root for MG);
%  finite mu                      : two-pressure closure;
%  mu = theta = inf, nu = 0       : pT closure;
%  finite theta                   : implicit heat exchange Delta Q(p);
%  mu = theta = inf, finite nu    : finite-rate mass transfer;
%  mu = theta = nu = inf          : pTg closure on the saturation curve;
% activation criteria (interface cells, T > T_sat) and projection of U^n.

\subsection{Asymptotic limits}
\label{sec:limits}
\begin{theorem}[Relaxation limits]\label{thm:limits}
% TODO (B7 Thms 1 and 2 merged):
% (i) rates -> infinity: O(1/(tau mu)), O(1/(tau theta)), O(1/(tau nu)) convergence to the
%     equilibrium closures; the limit scheme is a consistent discretisation of the
%     equilibrium model (Kapila volume equation);
% (ii) rates -> 0: the frozen scheme is recovered exactly.
\end{theorem}
\begin{proposition}[Low-Mach limit]\label{prop:lowmach}
% TODO (B7 Thm 3).
\end{proposition}

\subsection{Global nonlinear system and its solution}
\label{sec:solver}
% TODO: G_i(P) = \bar p^{cl}_i(P) - P_i with a 3-point (5-point) stencil;
% monotonicity -> diagonally dominant M-matrix;
% Newton with finite-difference Jacobian scaled by the natural pressure scale
% (residual scale / Jacobian scale); pressure-perturbation unknown P - P_ref;
% Thomas (1D) / BiCGStab with a two-level preconditioner (Jacobi + constant coarse space);
% convergence criterion and the round-off floor:
% (round-off floor of the closure, B9.5 Prop 11: implementation remark / appendix)

\subsection{Time integration}
\label{sec:time}
% TODO: semi-implicit IMEX Runge--Kutta (Boscarino--Filbet--Russo);
% LSDIRK2(2,2,2), gamma = 1 - 1/sqrt(2); explicit stage states;
% time step from the material speed only (material CFL);
% start-up ramp (first step at acoustic CFL 0.4, growth <= 1.2, first order while ramping);
% a-posteriori fallback (MOOD): an inadmissible stage -> same step at first order, then halve.

\subsection{Multi-dimensional extension}
\label{sec:multid}
% TODO: dimension-by-dimension faces; vertex-divergence form of D^*_f
% (reduces bitwise to 1D for y-uniform data); positivity bounds per direction (1/4);
% cross stencil for the shock budget.

\subsection{Summary of properties}
\label{sec:properties}
% TODO: short list with pointers: well-balanced contacts (R2), temperature
% consistency (R3), conservation (R1), mixture-energy consistency (R6),
% positivity (Theorem~\ref{thm:positivity}), limits (Theorem~\ref{thm:limits}).
% Remaining choices that are not derived: material CFL, local/global Rusanov
% speed, minmod, start-up ramp and MOOD constants, activation threshold eps_I.
